\section{Architektur Codetask}

Codetask besteht aus einer Microservice-Architektur, die grundlegend aus sieben Services besteht.

Das Backend besteht dabei aus drei Services. Dazu gehört die Haupt-API, aber auch die Check-API, in der ein Scala-Compiler läuft, um Code live auszuführen. In der sogenannten Codetask-Aufgaben können Anwender Code schreiben und prüfen, ob dieser korrekt ist.
Da Code live ausgeführt wird, wurde dieses Feature in der Check-API isoliert, damit die Haupt-API nicht angegriffen werden kann.
Die Haupt-API und die Check-API verwenden das Play3-Framework und auch Scala 3.
Der dritte Service für das Backend ist die AI-API. Diese wird genutzt, damit Studierende beim Bearbeiten von Aufgaben eine KI um Hilfe bitten können.
Aufgrund der für diese Implementierung benötigten Abhängigkeiten ist die AI-API in Scala 2 geschrieben und nutzt das Play2-Framework. Aus diesem Grund wurde sie von der Haupt-API isoliert.
Die Haupt-API und die AI-API sind auch von außen erreichbar und haben ihre eigene Host-DNS.
Alle drei APIs nutzen sbt als Build-Tool und da alle drei mit Scala geschrieben sind, sind diese auch stateless, da Scala eine funktionale Programmiersprache ist.

Für das Frontend werden zwei Services verwendet. Ein Service ist das Webframework, das aus dem Vue3-Framework besteht, der andere Service ist der dazugehörende Webserver nginx.

Die Datenbank ist ebenfalls ein eigener Service. Als Datenbanksystem wird die NoSQL-Technik MongoDB in der Version 6.0.0 verwendet, die das EOL erreicht hat.
Um die Verwendung von MongoDB zu vereinfachen, gibt es als Service auch Mongo-Express. Mit diesem kann die MongoDB-Anwendung über eine WebUI verwendet werden.

In der Produktions- und Staging-Umgebung gibt es zusätzlich den Service Traefik, der als Reverse Proxy verwendet wird.

Die Aufgaben von Codetask werden in eigenen GitLab-Repositories angelegt und gespeichert. Der Inhalt wird von der Haupt-API abgerufen und dann in der Datenbank gespeichert.
Dabei gibt es zwei verschiedene GitLab-Repos. Einmal ein öffentliches, in dem der Inhalt der Vorlesungen vorhanden ist, und ein privates, in dem Klausuren erstellt und gespeichert werden. 
Das private GitLab-Repository ist jedoch noch nicht im Code eingebunden. Damit die Haupt-API auf die Repos zugreifen kann, wird ein Token verwendet.

Alle GitLab-Repos liegen nicht auf der HTWG-Instance, sondern in der öffentlichen GitLab-Instance.